% Part: applied-modal-logics
% Chapter: epistemic-logic
% Section: language-epistemic-logic

\documentclass[../../../include/open-logic-section]{subfiles}

\begin{document}

\olfileid{aml}{el}{lan}

\olsection{Basa Logika Epistemik}

\begin{defn}
Upamane $G$ himpunan lambang agen. Basa dhasar logika epistemik
multiagen ngemot
\begin{enumerate}
  \tagitem{prvFalse}{Konstanta proposisional kanggo !!{falsity}~$\lfalse$.}{}
  \tagitem{prvTrue}{Konstanta proposisional kanggo !!{truth}~$\ltrue$.}{}
  \item Sawijining himpunan kang !!{denumerable}s saka !!{propositional variable}s: $\Obj
    p_0$, $\Obj p_1$, $\Obj p_2$, \dots
  \item Konektif proposisional: \startycommalist
  \iftag{prvNot}{\ycomma $\lnot$ (negasi)}{}%
  \iftag{prvAnd}{\ycomma $\land$ (konjungsi)}{}%
  \iftag{prvOr}{\ycomma $\lor$ (disjungsi)}{}%
  \iftag{prvIf}{\ycomma $\lif$ (!!{conditional})}{}%
  \iftag{prvIff}{\ycomma $\liff$ (!!{biconditional})}{}.
  \item Operator kawruh $\Knows_a$ kang $a \in G$.
\end{enumerate}
\end{defn}

Yen kita mung ngrembug kawruh siji agen ing sistem, rujukan marang
himpunan~$G$ lan agen-agen siji-siji bisa diilangi. Ing kasus iki,
kita mung nggunakake operator dhasar~$\Knows$.

\begin{defn}
\emph{!!^{formula}s} ing basa epistemik ditetepake kanthi induksi
  kaya mangkene:
\begin{enumerate}
\tagitem{prvFalse}{$\lfalse$ iku !!{formula} atomik.}{}

\tagitem{prvTrue}{$\ltrue$ iku !!{formula} atomik.}{}

\item Saben variabel proposisional $\Obj p_i$ iku !!{formula} (atomik).

\tagitem{prvNot}{Yen $!A$ iku !!a{formula}, $\lnot !A$ uga
  !!a{formula}.}{}

\tagitem{prvAnd}{Yen $!A$ lan $!B$ iku !!{formula}s, $(!A \land
  !B)$ uga !!a{formula}.}{}

\tagitem{prvOr}{Yen $!A$ lan $!B$ iku !!{formula}s, $(!A \lor !B)$
  uga !!a{formula}.}{}

\tagitem{prvIf}{Yen $!A$ lan $!B$ iku !!{formula}s, $(!A \lif !B)$
  uga !!a{formula}.}{}

\tagitem{prvIff}{Yen $!A$ lan $!B$ iku !!{formula}s, $(!A \liff !B)$
  uga !!a{formula}.}{}

\item Yen $!A$ iku !!a{formula} lan $a \in G$, $\Knows_a !A$
  uga !!a{formula}.

\tagitem{limitClause}{Ora ana liyane kang dadi !!a{formula}.}{}
\end{enumerate}
\end{defn}

Yen !!a{formula}~$!A$ ora ngemot $\Knows_a$ kanggo agen endi wae,
rumus iku diarani \emph{tanpa modalitas}.

\begin{defn}
  Operator $\Knows$ nglambangake kawruh individu, dene $\EKnows$,
  kang lumrahe diwaca ``saben wong ngerti,'' nglambangake kawruh
  klompok. Yen $G' \subseteq G$, kita netepake $\EKnows_{G'} !A$
  minangka cekakan kanggo \[\bigwedge_{b \in G'} \Knows_b !A.\]
\end{defn}

Kita uga bisa netepake pangertèn kawruh kang luwih kuwat, yaiku
\emph{kawruh umum} ing klompok agen~$G$. Yen sawijining informasi
dadi kawruh umum ing klompok agen, saben kombinasi agen ing klompok
iku ngerti yen liyane ngerti yen liyane ngerti \dots tanpa wates.
Iki luwih kuwat tinimbang kawruh klompok, lan gampang ditemokake
model relasional kang !!a{formula} dadi kawruh klompok nanging dudu
kawruh umum. Kita nganggo $\CKnows_G !A$ kanggo nglambangake
  ``ing~$G$, pratelan~$!A$ dadi kawruh umum.''

\end{document}
